\section{Random Forests}

The Python part of this exercise is available in \myhref{https://github.com/LosDelDGIIM/LosDelDGIIM.github.io/blob/main/subjects/Erasmus-DUE/GKI/Exercises/Exercise06.ipynb}{this Jupyter Notebook}.\\

\begin{ejercicio}
    What is the core principle behind Random Forests, and how does it differ from a single decision tree?

    The core principle behind Random Forests is the concept of ensemble learning, where multiple decision trees are trained on different subsets of the data and features. Each tree makes a prediction, and the final output is determined by aggregating the predictions (e.g., majority voting for classification or averaging for regression). This approach differs from a single decision tree, which relies on a single model that may be prone to overfitting and may not generalize well to unseen data. Random Forests, by combining multiple trees, can reduce overfitting and improve generalization performance.
\end{ejercicio}

\begin{ejercicio}
    What is Gini impurity?

    Gini impurity is a measure of the impurity or disorder in a dataset used in decision tree algorithms. It quantifies how often a randomly chosen element from the set would be incorrectly labeled if it were randomly labeled according to the distribution of labels in the subset. The formula for Gini impurity is:
    \[Gini = 1 - \sum_{i=1}^{n} p_i^2\]
    where \(p_i\) is the proportion of instances belonging to class \(i\) in the subset. A Gini impurity of 0 indicates that all instances belong to a single class, while a Gini impurity close to 1 indicates that the classes are evenly distributed.
\end{ejercicio}

\begin{ejercicio}
    Explain the concepts of ``bagging'' and ``feature randomness'' in the context of Random Forests.

    ``Bagging'' (Bootstrap Aggregating) is a technique used in Random Forests where multiple subsets of the training data are created by sampling with replacement. Each decision tree in the forest is trained on a different subset of the data, which helps to reduce variance and improve model robustness. 

    ``Feature randomness'' refers to the process of selecting a random subset of features for each split in the decision trees. This means that when a tree is being built, it does not consider all features for making splits but rather a random selection of them. This introduces additional diversity among the trees, which can further reduce overfitting and improve the overall performance of the Random Forest.
\end{ejercicio}

\begin{ejercicio}
    How do Random Forests handle classification and regression tasks differently?

    Random Forests handle classification and regression tasks by using different aggregation methods for the final predictions. In classification tasks, Random Forests use majority voting, where the class that receives the most votes from the individual trees is selected as the final prediction. In regression tasks, Random Forests use averaging, where the predicted values from all the trees are averaged to produce the final output. This allows Random Forests to effectively handle both types of tasks while maintaining their robustness and accuracy.
\end{ejercicio}

\begin{ejercicio}
    What are some key hyperparameters of Random Forests, and how do they impact model performance?

    Some key hyperparameters of Random Forests include:
    \begin{itemize}
        \item The number of trees in the forest. Increasing the number of trees can improve performance but also increases computational cost.
        \item The maximum depth of the trees. Limiting the depth can help prevent overfitting, while allowing deeper trees can capture more complex patterns in the data.
        \item The minimum number of samples required to split a node. Higher values can lead to simpler trees and reduce overfitting, while lower values can create more complex trees that may overfit the training data.
        \item The number of features to consider when looking for the best split. This can impact the diversity of the trees and, consequently, the overall performance of the model.
    \end{itemize}
\end{ejercicio}

\begin{ejercicio}
    Discuss the advantages and disadvantages of using Random Forests compared to other machine learning algorithms.

    Advantages of Random Forests include:
    \begin{itemize}
        \item They are robust to overfitting due to the ensemble nature of the model.
        \item They can handle both classification and regression tasks effectively.
        \item They can capture complex relationships in the data without requiring extensive feature engineering.
        \item They provide feature importance measures, which can help identify the most relevant features in the dataset.
    \end{itemize}

    Disadvantages of Random Forests include:
    \begin{itemize}
        \item They can be computationally expensive, especially with a large number of trees and features.
        \item They may not perform well on datasets with a large number of irrelevant features, as the random selection of features can lead to less informative splits.
        \item They can be less interpretable than simpler models like decision trees or linear regression, making it harder to understand the underlying relationships in the data.
    \end{itemize}
\end{ejercicio}